\section{术语与符号表}

本表按主题分七组列出正文用到的术语。每条先给英文与中文，再给一句话解释；标「待核」者表示笔者对经典文献中的精确表述尚无十足把握，读者引用前请核对原文。符号的基本约定见第 1 讲，此处不重复。

\begingroup
\small
\setlength{\tabcolsep}{4pt}
\begin{longtable}{@{}>{\raggedright\arraybackslash}p{0.285\textwidth}>{\raggedright\arraybackslash}p{0.66\textwidth}@{}}
\toprule
\textbf{英文（中文）} & \textbf{一句话解释} \\
\midrule
\endfirsthead
\toprule
\textbf{英文（中文）} & \textbf{一句话解释} \\
\midrule
\endhead
\bottomrule
\endlastfoot

\multicolumn{2}{@{}l}{\textbf{一、对偶、群论与耦合常数}}\\
\addlinespace[2pt]
\term{electric--magnetic duality}{电磁对偶} & 把电荷与磁荷、局域算子与 disorder 算子互换的对称性，同时把耦合常数取倒数。\\
\term{Montonen--Olive duality}{Montonen--Olive 对偶} & 1977 年提出的猜想：$N=4$ 超对称规范理论在强弱耦合互换下自对偶。\\
\term{S-duality}{S-对偶} & $\tau\mapsto-1/(n_g\tau)$ 的强弱对偶变换，交换 Wilson 算子与 't Hooft 算子。\\
\term{T-duality}{T-对偶} & 环面上半径 $R\mapsto 1/R$ 的等价性；在 Hitchin 纤维上表现为对偶环面。\\
\term{GNO duality}{GNO 对偶} & Goddard--Nuyts--Olive 的结论：电、磁荷取值格的对偶确定对偶群。\\
\term{Langlands dual group}{Langlands 对偶群 $\Ldual G$} & 把根系与余根系互换所得的复群，例如 $SU(N)$ 与 $PSU(N)$ 互换。\\
\term{lacing number}{织数 $n_g$} & 描述 Dynkin 图边长不对称的整数：单边型为 $1$，$B_n$、$C_n$、$F_4$ 为 $2$，$G_2$ 为 $3$。\\
\term{Hecke group}{Hecke 群} & $G_2$、$F_4$ 时 $S$ 生成的 $SL(2,\RR)$ 的无限离散子群，代替 $SL(2,\ZZ)$。\\
\term{duality group}{对偶群} & 由 $S$ 与 $T$ 生成的离散群，单边型时为 $SL(2,\ZZ)$。\\
\term{complexified coupling}{复耦合 $\tau$} & $\tau=\theta/2\pi+4\pi\ii/e^2$，实部与虚部分别编码 $\theta$ 与 $1/e^2$。\\
\term{theta angle}{$\theta$ 角} & 瞬子数项的系数除以 $2\pi$；$\theta\mapsto\theta+2\pi$ 是规范等价，故 $\Re\tau$ 周期为 $1$。\\
\term{instanton number}{瞬子数} & $(1/8\pi^2)\int\Tr F\wedge F$ 的整数值，标记规范场的拓扑类。\\
\term{weight lattice}{权格 $P$} & 表示权所在的格；$SU(2)$ 的权格为 $(1/2)\ZZ$，$SO(3)$ 的为 $\ZZ$。\\
\term{coweight lattice}{余权格 $P^\vee$} & 磁荷所在的格，与权格互为对偶，GNO 对偶据此定义。\\

\addlinespace[3pt]
\multicolumn{2}{@{}l}{\textbf{二、超对称与扭化}}\\
\addlinespace[2pt]
\term{$\mathcal N=4$ super Yang--Mills}{$N=4$ 超对称 Yang--Mills} & 四维最大超对称规范理论，含 $16$ 个超荷，共形且具有 S-对偶。\\
\term{supercharge}{超荷} & 超对称变换的生成元，把玻色子变为费米子、费米子变为玻色子。\\
\term{topological supercharge}{拓扑超荷} & 扭化后成为 BRST 算子的标量超荷，满足 $Q^2=0$。\\
\term{twisting}{扭化} & 把 R-对称与时空旋转群对角化，从超对称理论得到拓扑场论的操作。\\
\term{GL twist}{GL 扭化} & Kapustin--Witten 所用的特定扭化，因与几何 Langlands 纲领相关而得名。\\
\term{canonical parameter}{典范参数 $\Psi$} & 扭化参数 $t$ 与 $\tau$ 的组合 $\Psi=(\tau t^2+\bar\tau)/(1+t^2)$，是 S-对偶下的自然耦合。\\
\term{twisting parameter}{扭化参数 $t$} & 扭化族的 $\PP^1$ 参数；$t$ 与 $-t$ 等价，$t=\pm\ii$ 对应 $\Psi=\infty$。\\
\term{BRST operator}{BRST 算子} & 扭化理论中幂零的标量超荷 $Q$，物理态由其上同调给出。\\
\term{$Q$-cohomology}{$Q$-上同调} & 以 BRST 微分取上同调，得到拓扑不变的态空间与算子代数。\\
\term{$R$-symmetry}{R-对称} & 在超代数中旋转超荷的对称群；$N=4$ 时为 $SU(4)\cong Spin(6)$。\\

\addlinespace[3pt]
\multicolumn{2}{@{}l}{\textbf{三、Hitchin 模空间与超凯勒几何}}\\
\addlinespace[2pt]
\term{Riemann surface}{Riemann 面 $C$} & 紧复曲线；四维理论沿 $C$ 紧化，$C$ 是二维理论的底空间。\\
\term{genus}{亏格 $g$} & 复曲线的洞数；$\MH$ 的维数与 Hitchin 基的维数都随 $g$ 变化。\\
\term{$G$-bundle}{$G$-丛} & 以群 $G$ 为结构群的纤维丛；其模空间记为 $\mathrm{Bun}_G$。\\
\term{principal bundle}{主丛} & 纤维与 $G$ 同构且带 $G$ 右作用的丛，是 Higgs 丛的底。\\
\term{Higgs bundle}{Higgs 丛 $(E,\varphi)$} & 全纯 $G$-丛 $E$ 配一个全纯的 $\mathrm{ad}$-值 $(1,0)$ 形式 $\varphi$。\\
\term{Higgs field}{Higgs 场 $\varphi$} & $\varphi\in H^0(C,\mathrm{ad}(E)\otimes K_C)$，即 $\mathrm{ad}(E)$ 值全纯 $(1,0)$ 形式。\\
\term{Hitchin equations}{Hitchin 方程} & $F_A+[\varphi,\varphi^*]=0$ 与 $\bar\partial_A\varphi=0$，是 $\MH$ 上的自对偶方程。\\
\term{Hitchin moduli space}{Hitchin 模空间 $\MH$} & Higgs 丛（等价地 Hitchin 方程解）的模空间，是超凯勒流形。\\
\term{stability}{稳定性} & 排除破坏模空间良好性的子丛，Higgs 丛的稳定性使模空间为流形。\\
\term{hyperkähler}{超凯勒} & 带有满足四元数关系的三个可积复结构的度量几何；$\MH$ 是超凯勒流形。\\
\term{hyperkähler quotient}{超凯勒商} & 三个矩映射条件联立后商去群作用，$\MH$ 即无穷维超凯勒商。\\
\term{complex structures $I,J,K$}{复结构 $I,J,K$} & 三个满足 $I^2=J^2=K^2=IJK=-1$ 的可积复结构，膜类型按它们区分。\\
\term{symplectic form}{辛形式} & 闭非退化 $2$-形式；$\omega_J$、$\omega_K$ 恰当而 $\omega_I$ 给出拓扑类。\\
\term{moment map}{矩映射 $\mu$} & 哈密顿群作用的生成函数；辛商取 $\mu^{-1}(0)$ 再商群作用。\\
\term{Hitchin fibration}{Hitchin 纤维化} & 由 Casimir 不变量给出的映射 $\MH\to B$，其纤维是复环面。\\
\term{completely integrable}{完全可积} & Hitchin 系统：存在 $\dim B$ 个泊松对合的函数，其共同水平集为环面。\\
\term{complex torus}{复环面} & 形如 $\CC^n/\Lambda$ 的紧复流形；$\MH$ 的 generic 纤维就是复环面。\\
\term{spectral curve}{谱曲线} & 由 $\det(\lambda-\varphi)=0$ 定义的曲线，用来描述 Hitchin 纤维。\\

\addlinespace[3pt]
\multicolumn{2}{@{}l}{\textbf{四、sigma 模型、广义复几何与膜}}\\
\addlinespace[2pt]
\term{sigma model}{sigma 模型} & 二维场论，基本场是从世界面到靶空间的映射，作用量含靶空间度量。\\
\term{A-model}{A-模型} & 依赖辛结构、只保留 $(1,1)$ 类超荷的拓扑 sigma 模型。\\
\term{B-model}{B-模型} & 依赖复结构、保留全纯超荷的拓扑 sigma 模型。\\
\term{mirror symmetry}{镜像对称} & 两个不同靶空间上 A-模型与 B-模型等价；本文中靶空间为同一 $\MH$ 的不同复结构。\\
\term{generalized complex geometry}{广义复几何} & 统一复结构与辛结构的框架，用 $T\oplus T^*$ 上的结构描述膜的类型。\\
\term{brane}{膜} & sigma 模型中的边界条件；在拓扑场论中也指可被线算子作用的对象。\\
\term{A-brane}{A-膜} & A-模型允许的边界条件：Lagrange 或余各向同性子流形加平坦丛。\\
\term{B-brane}{B-膜} & B-模型允许的边界条件：复子流形加全纯 Chan--Paton 丛。\\
\term{Lagrangian submanifold}{Lagrange 子流形} & 辛形式在其上限制为零且维数为靶空间一半的子流形。\\
\term{coisotropic submanifold}{余各向同性子流形} & 法向方向上辛形式退化（含 Lagrange 情形为特例）的子流形。\\
\term{Chan--Paton bundle}{Chan--Paton 丛} & 膜上的向量丛，编码开弦端点的内部自由度。\\
\term{SYZ conjecture}{SYZ 猜想} & 镜像流形由互为对偶的环面纤维化相联系（Strominger--Yau--Zaslow）。\\

\addlinespace[3pt]
\multicolumn{2}{@{}l}{\textbf{五、线算子与特征膜}}\\
\addlinespace[2pt]
\term{Wilson operator}{Wilson 算子} & 沿闭曲线的和乐之迹 $W_R(S)=\Tr_R P\exp(-\oint A)$，是电型线算子。\\
\term{'t Hooft operator}{'t Hooft 算子} & 先规定沿曲线的奇异性再作路径积分的 disorder 算子，带磁荷标签。\\
\term{line operator}{线算子} & 支撑在一维曲线上的算子；在二维理论中作用在膜上。\\
\term{loop operator}{圈算子} & 支撑在闭曲线上的线算子，其周边界与膜相交。\\
\term{order operator}{有序算子} & 由局域场（如联络）构造的算子，Wilson 算子是其典型代表。\\
\term{disorder operator}{无序算子} & 由奇异性规定而非局域场构造的算子，'t Hooft 算子是其代表。\\
\term{eigenbrane}{特征膜} & 满足 $XB=B\otimes V$ 的膜，即线算子的本征对象；对应几何 Langlands 的特征层。\\
\term{electric eigenbrane}{电特征膜} & Wilson 算子族的联合特征膜，自动带一个平坦向量丛。\\
\term{magnetic eigenbrane}{磁特征膜} & 't Hooft 算子族的联合特征膜，即 Hitchin 纤维上的 type $F$ 膜。\\
\term{zerobrane}{零膜（点膜）} & 支撑在单点上、Chan--Paton 为天穹层的 $(B,B,B)$ 型膜。\\
\term{local system}{局部系统} & 基本群 $\pi_1(C)$ 的表示，等价于局部常值层或平坦联络。\\
\term{fundamental group}{基本群 $\pi_1(C)$} & 曲线上闭路同伦类的群；$\Ldual G$-局部系统即其到 $\Ldual G_\CC$ 的同态。\\
\term{brane type}{膜类型} & 对 $I,J,K$ 三个复结构分别标 $A$ 或 $B$ 所得的三字母标签。\\

\addlinespace[3pt]
\multicolumn{2}{@{}l}{\textbf{六、Hecke 修改、单极子与仿射 Grassmannian}}\\
\addlinespace[2pt]
\term{Hecke operator}{Hecke 算子} & 几何 Langlands 中把丛在指定点处改写后拉回、再推前的对应算子。\\
\term{Hecke modification}{Hecke 修改} & 在点 $p_0$ 外保持同构、在 $p_0$ 处允许极点的丛的重写。\\
\term{Hecke correspondence}{Hecke 对应} & $\mathrm{Bun}\times\mathrm{Bun}$ 中记录修改关系的子簇，由此定义 Hecke 算子。\\
\term{Hecke eigensheaf}{Hecke 特征层} & 在 Hecke 算子下被送到自身张量一个向量空间的层，对应特征膜。\\
\term{Bogomolny equations}{Bogomolny 方程} & 四维自对偶方程降维所得的单极子方程，形如 $F=\ii\star D\varphi$。\\
\term{monopole}{单极子} & 三维空间带磁荷的有限能量解；'t Hooft 算子的奇异性对应它。\\
\term{monopole bubbling}{单极子冒泡} & 单极子奇点处瞬子尺度退化为零、并入单极子的经典现象。\\
\term{affine Grassmannian}{仿射 Grassmannian} & 圈群商 $G_\CC(X)/G_\CC(O)$，参数化在一点处的 Hecke 修改。\\
\term{Schubert cell}{Schubert 胞腔} & 仿射 Grassmannian 中按支配余权的 $G_\CC(O)$-轨道分层所得胞腔。\\
\term{Schubert cycle}{Schubert 循环} & Schubert 胞腔的闭包，出现在 Hecke 对应的纤维中。\\

\addlinespace[3pt]
\multicolumn{2}{@{}l}{\textbf{七、几何 Langlands、$D$-模与量子化}}\\
\addlinespace[2pt]
\term{geometric Langlands program}{几何 Langlands 纲领} & $\Ldual G$-局部系统与 $D$-模（或特征层）之间的对应，是本文的主题。\\
\term{Langlands program}{Langlands 纲领} & 数论与表示论中的巨大对偶网；几何版本用曲线上的丛替换数域。\\
\term{eigensheaf}{特征层} & Hecke 算子的本征对象；KW 在拓扑场论一侧把对应物称为特征膜。\\
\term{$D$-module}{$D$-模} & 微分算子层 $D$ 上的模，几何 Langlands 中对应局部系统的一侧。\\
\term{sheaf of differential operators}{微分算子层 $D$} & 由向量场及其积生成的层；c.c.\ 开弦代数取值于其上。\\
\term{twisted $D$-module}{扭化 $D$-模} & 系数扭化或带中心特征的 $D$-模（待核：不同文献的扭化方式略有差异）。\\
\term{c.c.\ A-brane}{c.c.\ 膜} & 秩 $1$、满足 $(\omega^{-1}F)^2=-1$ 的典范余各向同性 A-膜，型为 $(A,B,A)$。\\
\term{open string BRST}{开弦 BRST} & 边界条件所保留的超荷在边界上成为 BRST 微分，其幂零性给出弦上同调。\\
\term{quantization parameter}{量子化参数 $f(\Psi)$} & c.c.\ 弦代数是作用在 $K_M^{1/2}\otimes L^{f(\Psi)}$ 上的微分算子层，$f(\Psi)=\Psi$。\\
\term{three brane types}{三类膜} & 实际出现的 $(B,B,B)$、$(B,A,A)$、$(A,B,A)$ 三种膜；$(B,B,A)$ 与 $(A,A,A)$ 不出现。\\
\term{ramification}{分歧} & 局部系统在一点处有非平凡单值性；几何 Langlands 向分歧情形推广。\\
\term{unramified}{非分歧} & 局部系统在每点局部平凡，即经典（非分歧）几何 Langlands 的情形。\\
\term{wild ramification}{野分歧} & 局部单值性含幂零部分、需额外数据刻画的情形（待核）。\\
\end{longtable}
\endgroup
